\documentclass[11pt]{article}
\usepackage{amsmath,amssymb,amsthm,geometry,hyperref}
\geometry{margin=1in}
\newtheorem{proposition}{Proposition}
\newtheorem{lemma}{Lemma}
\newcommand{\R}{\mathbb R}
\newcommand{\Id}{\operatorname{Id}}
\title{Exact first-stage steering audit with the original base scales retained}
\author{}
\date{8 September 2026}
\begin{document}
\maketitle

This independent audit concerns the first stage in Alp\"oge--Buckmaster,
\emph{Blowup for the Boussinesq equations with smooth forcing}, released
76-page PDF: Lemma 1.2, pp.~6--7; equations (3.2)--(3.5), pp.~8--9;
Lemmas 3.7--3.8 and equations (3.24)--(3.34), pp.~15--18; equation
(3.58), p.~27; and the holding calculation, pp.~29 and 35--36.
The downloaded source is
\texttt{research/alpoge\_buckmaster\_boussinesq.pdf} in the parent notebook.
All objects below retain their laboratory coordinates and gravity direction.
The affine wave is local; this audit asserts no compact spatial support for it.

\paragraph{Parameters and physical scaling.}
Let $A_0>0$, $\sigma _0=\nu_{\rm time}=\sqrt{A_0}$, and
$\mu_{\rm sp}=\lambda _0>0$. The original datum of the paper additionally
has $A_0\geq1$ and $\lambda _0=\lceil\sqrt{A_0}/2\rceil$.
Let $\lambda _1>0$ be the paper's first frequency and put
$\Lambda=\mu_{\rm sp}\lambda _1$. The time scaling parameter
$\nu_{\rm time}$ is not physical viscosity. Physical viscosity will be
denoted $\nu_{\rm phys}$ and scalar diffusivity $\kappa_{\rm phys}$.
The paper's first pulse compression is denoted $C=\Lambda _1$ here;
it is not the physical frequency $\Lambda$. Its pulse mass parameter,
called $\mu$ in (3.25), is denoted $m$ here and is not $\mu_{\rm sp}$.

For $(X,T)=(\mu_{\rm sp}x,\nu_{\rm time}t)$ the exact physical map is
\begin{align}
 \theta(x,t)&=\frac{\nu_{\rm time}^2}{\mu_{\rm sp}}
                   \widetilde\theta(X,T),&
 u(x,t)&=\frac{\nu_{\rm time}}{\mu_{\rm sp}}
                   \widetilde u(X,T),\\
 p(x,t)&=\frac{\nu_{\rm time}^2}{\mu_{\rm sp}^2}
                   \widetilde p(X,T),&
 \omega(x,t)&=\nu_{\rm time}\widetilde\omega(X,T).
\end{align}
In particular $\Theta=(A_0/\mu_{\rm sp})\widetilde\Theta$,
$\Omega=\sigma _0\widetilde\Omega$,
$D=\sigma _0\widetilde D$, $G=A_0\widetilde G$, and
$\zeta(t)=\widetilde\zeta(\sigma _0t)$. These statements follow by
differentiating each displayed expression: spatial derivatives contribute
$\mu_{\rm sp}$ and time derivatives contribute $\sigma _0$.
For example $\partial_t\Theta=(\sigma _0^3/\mu_{\rm sp})
\partial_T\widetilde\Theta$, and
\[
-\frac{J\zeta\cdot G}{\Lambda|\zeta|^2}\Omega
=\frac{\sigma _0^3}{\mu_{\rm sp}}
 \left(-\frac{J\widetilde\zeta\cdot\widetilde G}
 {\lambda _1|\widetilde\zeta|^2}\widetilde\Omega\right).
\]
Likewise $\partial_t\Omega=\sigma _0^2\partial_T\widetilde\Omega
=\Lambda\zeta_1\Theta$. This proves the amplitude-system morphism.
For retained diffusion the corresponding coefficients are
$\widetilde\nu_{\rm phys}=\nu_{\rm phys}\mu_{\rm sp}^2/\sigma _0$
and $\widetilde\kappa_{\rm phys}=\kappa_{\rm phys}\mu_{\rm sp}^2/\sigma _0$,
because the Laplacian contributes $\mu_{\rm sp}^2$.

\paragraph{Rigid geometry and its exact equations.}
Set
\[
J=\begin{pmatrix}0&-1\\1&0\end{pmatrix},\quad
R_\alpha=\begin{pmatrix}\cos\alpha&-\sin\alpha\\
\sin\alpha&\cos\alpha\end{pmatrix},\quad
e(s)=(\sin s,\cos s).
\]
Fix $s$ independent of time. During the first stage put
\begin{align}
 D&=\dot\alpha J,& M&=R_\alpha,&
 G&=-A_0R_\alpha e_2=(A_0\sin\alpha,-A_0\cos\alpha),\\
 \zeta&=R_\alpha e(s)=(\sin(s-\alpha),\cos(s-\alpha)).
\end{align}
These formulas have the exact derivative equations
\begin{align}
 \dot M&=DM,& \det M&=1,& M^{-T}e(s)&=\zeta,\\
 \dot G&=\dot\alpha JG=-D^TG,&
 \dot\zeta&=\dot\alpha J\zeta=-D^T\zeta,\\
 \dot\zeta_1&=-\dot\alpha\zeta_2,&
 \dot\zeta_2&=\dot\alpha\zeta_1,& |\zeta|&=1.
\end{align}
For verification, $R_\alpha'=JR_\alpha$, $J^T=-J$, and direct
matrix multiplication gives all the formulas. Also
\[
 J\zeta\cdot G
=-A_0\cos(s-\alpha)\sin\alpha
 -A_0\sin(s-\alpha)\cos\alpha=-A_0\sin s.
\]
Consequently the physical first-layer equations are exactly
\begin{equation}\label{eq:physical}
 \dot\Theta=a\Omega,\qquad\dot\Omega=b\Theta,\qquad
 a=\frac{A_0\sin s}{\Lambda},\qquad b=\Lambda\zeta_1.
\end{equation}
The base gradient evolves by the displayed transport law even though its
norm is constant. No additional older shear occurs at the first stage.
For the affine background $u_b=Dx$, $\theta_b=G\cdot x$, an exact
pressure and force are
\[
 p_b=\frac{\dot\alpha^2}{2}|x|^2+\frac12(G\cdot x)x_2,
 \qquad f_{u,b}=\left(\ddot\alpha-\frac{G_1}{2}\right)Jx,
 \qquad f_{\theta,b}=0.
\]
Indeed $(\partial_t+u_b\cdot\nabla)\theta_b
=(\dot G+D^TG)\cdot x=0$ and
$\partial_tu_b+u_b\cdot\nabla u_b=(\ddot\alpha J-\dot\alpha^2\Id)x$.
The gradient of the displayed pressure cancels the quadratic-rotation
term and leaves
\[
 (\ddot\alpha J)x+\tfrac12[Gx_2-e_2(G\cdot x)]
 =\left(\ddot\alpha-\frac{G_1}{2}\right)Jx
\]
after subtracting $\theta_be_2$. This proves the momentum equation
and force sign in the original laboratory coordinates.
Both affine Laplacians vanish, so these background formulas also hold
with either retained diffusion coefficient. During first holding the
background force is $-(A_0\sin s/2)Jx$, whose curl is
$-A_0\sin s$; it does not vanish merely because $D=0$.

\paragraph{The full pulse family.}
Choose $h\in C_c^\infty((0,1))$ with $h\geq0$ and
$\int_0^1h=1$, let $H=\|h\|_\infty$, and choose
\[
 c_p>1,\qquad C\geq2c_pH,\qquad
 0<s\leq\tfrac18,\qquad c_pCH\sin s\leq\tfrac14.
\]
The paper takes $c_p=3$, $h=\eta_2$, and $C=\Lambda _1$.
Let $\eta$ be a smooth nondecreasing step, zero on $(-\infty,0]$
and one on $[1,\infty)$. For $m\in[0,c_p]$ set
\[
 \tau_b=1+C^{-1},\qquad
 z_m(\tau)=\begin{cases}1-\eta(\tau),&0\leq\tau\leq1,\\
 -mC h(C(\tau-1)),&1\leq\tau\leq\tau_b.
 \end{cases}
\]
The flatness of the step and the compact interior support of $h$ make
this join smooth. Let $P_m$ solve the linear initial value problem
\begin{equation}\label{eq:P}
 P_m''=z_mP_m,\qquad P_m(0)=P_m'(0)=1.
\end{equation}
Linear ODE existence applies on the entire finite interval independently
of any quotient. The next proof verifies positivity before defining a
quotient globally.

\begin{proposition}\label{prop:shoot}
For every $m\in[0,c_p]$, $P_m>0$ on $[0,\tau_b]$. With
$v=P_m'/P_m$ and $d=v(1)$, one has
\begin{align}
 &(1+\tau)^{-1}\leq v(\tau)\leq1\quad(0\leq\tau\leq1),
 \qquad \log2\leq\log P_m(1)\leq1,\qquad \tfrac12\leq d\leq1,
 \\
 &\tfrac12\leq\chi(y):=\frac{P_m(1+y/C)}{P_m(1)}
 \leq1+y/C,\qquad -4c_p\leq V(y):=v(1+y/C)\leq d,
 \quad0\leq y\leq1,\\
 &0\leq\log P_m(\tau)\leq1+C^{-1}\quad(0\leq\tau\leq\tau_b).
\end{align}
The map $E(m)=v(\tau_b;m)$ is smooth and strictly decreasing, with
\begin{equation}\label{eq:rootderivative}
 E'(m)=-\int_0^1 h(y)
 \exp\!\left(-\frac2C\int_y^1V(r;m)\,dr\right)dy<0.
\end{equation}
It has a unique zero $m_*$, and
\[
 \tfrac12-\frac1{4C}\leq m_*\leq1,\qquad
 0\leq V\leq d\text{ at }m=m_*,\qquad
 \log2\leq\log P_{m_*}(\tau_b)\leq1+C^{-1}.
\]
Uniformly over all trials,
$-e^{8c_p/C}\leq E'(m)\leq-e^{-2/C}$.
At the selected root the sharper bounds are
$-1\leq E'(m_*)\leq-e^{-2/C}$.
\end{proposition}
\begin{proof}
On $[0,1]$, while $P>0$, its equation gives $P''\geq0$,
hence $P'\geq1$ and $P\geq1+\tau$. These inequalities exclude
a first zero by continuity. The quotient is therefore defined and obeys
$v'=1-\eta-v^2$. The constant $1$ is a supersolution.
The positive function $w=(1+\tau)^{-1}$ satisfies $w'=-w^2$.
Writing the differential equations for $v-1$ and $v-w$, and multiplying
by the respective positive integrating factors for their linear terms,
proves $w\leq v\leq1$. Integration proves the first displayed bounds.
The data $P(1),d$ do not depend on $m$.

On the pulse $\chi''+(mh/C)\chi=0$, $\chi(0)=1$,
$\chi'(0)=d/C$. Before a possible first zero, concavity gives
$\chi\leq1+dy/C\leq1+y/C<2$, since $C>2$.
The condition $mh/C\leq1/2$ gives $\chi''\geq-1$, hence
$\chi(y)\geq1+dy/C-y^2/2\geq1/2$. This excludes the zero and
proves the bound on the full pulse. The exact integral identity
\[
 \chi'(y)=\frac dC-\frac mC\int_0^yh(r)\chi(r)\,dr
 \geq-\frac{2c_p}C
\]
now gives $V=C\chi'/\chi\geq-4c_p$. Differentiating this quotient
gives the exact equation
\begin{equation}\label{eq:V}
 \frac{dV}{dy}=-mh(y)-\frac{V^2}{C}\leq0,\qquad V(0)=d,
\end{equation}
so $V\leq d\leq1$. Positivity and smooth parameter dependence of the
linear equation give smoothness of the quotient in $m$. Its derivative
$W=\partial_mV$ solves $W'+(2V/C)W=-h$, $W(0)=0$.
The integrating factor gives \eqref{eq:rootderivative}. The integrand is
positive wherever $h>0$, which is a nonempty set of positive measure.
The all-trial bounds on $V$ give the two exponential derivative bounds.

For $m=0$, solving \eqref{eq:V} gives
$E(0)=d/(1+d/C)>0$. For $m=c_p$, integration gives
$E(c_p)=d-c_p-C^{-1}\int_0^1V^2\leq1-c_p<0$.
Continuity and strict monotonicity give exactly one root in $(0,c_p)$.
There $V$ decreases from $d$ to zero, so $0\leq V\leq d$, and
\[
 m_*=d-\frac1C\int_0^1V^2\,dy\in[d-d^2/C,d].
\]
The derivative of $x-x^2/C$ is $1-2x/C>0$ on $[1/2,1]$.
This proves the stated interval for $m_*$. The selected derivative
bounds follow from $0\leq V\leq1$ in \eqref{eq:rootderivative}.
For every trial, on the pulse,
\[
 \log P_m(1+y/C)=\log P_m(1)+\log\chi(y)
 \in[\log2-\log2,\,1+\log(1+C^{-1})]
 \subset[0,1+C^{-1}].
\]
At the selected root the pulse contribution is
$C^{-1}\int_0^1V\,dy\in[0,C^{-1}]$; adding the first-part
integral proves the final gain. This also proves all-trial amplitude
bounds without imposing all-trial positivity of $V$.
\end{proof}

If the independent first-part exit datum $d$ is varied near its actual
value, differentiating \eqref{eq:V} gives
\[
 \partial_dE=\exp\!\left(-\frac2C\int_0^1V(r)dr\right),\qquad
 \frac{dm_*}{dd}=
 \frac{\exp[-(2/C)\int_0^1V]}{
 \int_0^1h(y)\exp[-(2/C)\int_y^1V]dy}>0.
\]
The implicit function theorem applies because $E'(m_*)<0$.
At the actual selected root the last derivative lies in
$[e^{-2/C},e^{2/C}]$. This derivative statement concerns the displayed
family with $d$ variable; it introduces no change to the paper's actual
datum $d$.

\paragraph{Control map, amplitudes, and every join.}
Set $\Gamma=\sigma _0\sin s$, $\tau=\Gamma(t-t_1)$,
$Z(t)=\sin s\,z_m(\tau)$, and
\[
 \alpha(t)=s-\arcsin Z(t).
\]
Then $\zeta_1=Z$, $\zeta_2=\sqrt{1-Z^2}\geq\sqrt{15}/4$,
and $0\leq\alpha<1/2$. Indeed $|Z|\leq c_pCH\sin s\leq1/4$,
$Z\leq\sin s$, and the increasing principal arcsine satisfies
$\arcsin Z\leq s$ and $\arcsin(1/4)\leq1/3$.
Differentiation yields the physical feedback exactly:
\[
 \dot\alpha=-\frac{\dot Z}{\sqrt{1-Z^2}}
            =-\frac{\dot Z}{\zeta_2},\qquad
 \dot G=\dot\alpha JG,\qquad \dot\zeta=\dot\alpha J\zeta.
\]
This is the paper's control with $F_1=0$. Given $P_{\rm ent}>0$, set
\begin{equation}\label{eq:amplitudes}
 \Theta(t)=-P_{\rm ent}P_m(\tau),\qquad
 \Omega(t)=-\frac\Lambda{\sigma _0}P_{\rm ent}P_m'(\tau).
\end{equation}
Their initial values are $\Theta(t_1)=-P_{\rm ent}$ and
$\Omega(t_1)=(\Lambda/\sigma _0)\Theta(t_1)$.
Since $\Gamma\Lambda/\sigma _0=\Lambda\sin s$ and
$a\Lambda/\sigma _0=\Gamma$, differentiation of
\eqref{eq:amplitudes} proves both equations \eqref{eq:physical}.
It also proves, without changing the original amplitudes,
\[
 v=\frac{\sigma _0\Omega}{\Lambda\Theta},\qquad
 \partial_t\log(-\Theta)=\Gamma v,\qquad
 |\Omega|=\frac\Lambda{\sigma _0}|v|(-\Theta).
\]
At the selected root $\Omega(t_b)=0$. Differentiating its endpoint
map in $m$ gives the explicit nonzero derivative
\[
 \left.\partial_m\Omega(t_b;m)\right|_{m=m_*}
=-\frac\Lambda{\sigma _0}P_{\rm ent}P_{m_*}(\tau_b)E'(m_*)>0.
\]
Near $\tau_b$, $h$ is identically zero, so $P''=0$.
The selected endpoint condition $P'(\tau_b)=0$ forces $P'=0$
throughout that neighborhood. Thus $Z=0$, $\alpha=s$, $\Omega=0$
and $\Theta$ is constant there. Extending these constants gives a smooth
stationary continuation and vanishing of every positive-order time
derivative of $\alpha,\Theta,\Omega$ at the join. Near $\tau=0$,
all derivatives of $\eta$ vanish to infinite order, so the pulse law
joins the preceding constant-$Z$ law smoothly as well.

For an explicit control bound at each derivative order $n\geq1$,
\[
 \|\partial_t^n Z\|_\infty\leq
 (\sin s)\Gamma^n
 \max\{\|\eta^{(n)}\|_\infty,c_pC^{n+1}\|h^{(n)}\|_\infty\}.
\]
The maximum is valid because the two pieces have disjoint interiors.
In particular
\[
 \|\dot\alpha\|_\infty\leq\frac4{\sqrt{15}}\Gamma\sin s
 \max\{\|\eta'\|_\infty,c_pC^2\|h'\|_\infty\},\qquad
 \int_{t_1}^{t_b}|\dot\alpha|dt\leq
 \frac4{\sqrt{15}}\sin s\left(1+c_pC\int_0^1|h'|\right).
\]
The first inequality follows from the feedback and the lower cosine
bound. For the second, integrate $|\dot Z|$: the step has variation
$\sin s$ and the pulse variation is at most
$c_pC\sin s\int|h'|$. All higher angle derivatives are explicit:
\[
 \partial_t^n\alpha=-\sum_{k=1}^n
 \frac{R_{k-1}(Z)}{(1-Z^2)^{k-1/2}}
 B_{n,k}(Z',\ldots,Z^{(n-k+1)}),
\]
where $R_0=1$, $R_j=(1-z^2)R'_{j-1}+(2j-1)zR_{j-1}$, and
\[
 B_{n,k}(x_1,\ldots,x_{n-k+1})=
 \sum_{\substack{\sum j\ell_j=n\\\sum\ell_j=k}}
 \frac{n!}{\prod_j\ell_j!}\prod_j\left(\frac{x_j}{j!}\right)^{\ell_j}.
\]
Repeated differentiation of $(\arcsin z)'=(1-z^2)^{-1/2}$ proves
the recurrence; repeated product and chain rules give the displayed
finite sum. The preceding bounds and $1-Z^2\geq15/16$ quantify each
derivative without an unspecified control singularity.

\paragraph{Growth, transition, and their exact physical gain.}
Let $k_1\geq0$ and $L_1=(k_1+41/8)\log\lambda _1>0$.
The physical seed, threshold, and initial time dictated by Lemma 1.2 are
\[
 P_{\rm seed}=\frac{A_0}{\mu_{\rm sp}}\lambda _1^{-k_1-6},\quad
 P_a=\frac{A_0}{\mu_{\rm sp}}\lambda _1^{-7/8},\quad
 t_{\rm in}=\sigma _0^{-1}.
\]
With $\alpha=0$ and $\Omega=(\Lambda/\sigma _0)\Theta$,
\eqref{eq:physical} gives
$\Theta=-P_{\rm seed}e^{\Gamma(t-t_{\rm in})}$.
Because $P_{\rm seed}e^{L_1}=P_a$, the first threshold occurs exactly at
$t_a=t_{\rm in}+L_1/\Gamma$, and its derivative is
$\partial_t(-\Theta)(t_a)=\Gamma P_a>0$.
The first transition has both feedbacks zero and length $\Gamma^{-1}$;
therefore $t_1=t_a+\Gamma^{-1}$ and $P_{\rm ent}=eP_a$.
Its contribution to the logarithmic gain is exactly one.
At the selected endpoint
\begin{align}
 t_b-t_{\rm in}&=\frac{L_1+2+C^{-1}}\Gamma,\\
 1+\log2&\leq\log\frac{P_f}{P_a}\leq2+C^{-1},
 \qquad P_f=-\Theta(t_b)=eP_aP_{m_*}(\tau_b),\\
 2eA_0\lambda _1^{1/8}&\leq A_1:=\Lambda P_f
 \leq e^{2+C^{-1}}A_0\lambda _1^{1/8}.
\end{align}
The gain for every trial through steering is between $1$ and
$2+C^{-1}$ relative to $P_a$, directly from Proposition~\ref{prop:shoot}.
The selected gain from the seed is between
$L_1+1+\log2$ and $L_1+2+C^{-1}$.

\paragraph{The actual endpoint background and first holding interval.}
At the selected endpoint the new layer is fully activated,
$\alpha=s$, $\zeta=e_2$, and $\Omega=0$. Its origin-gradient
increment is $\Lambda\Theta\zeta=-A_1e_2$. Thus the next
background gradient is exactly
\begin{equation}\label{eq:nextG}
 G_{<2}=G+\Lambda\Theta\zeta
 =(A_0\sin s,-A_0\cos s-A_1),\qquad
 \sigma _1^2=|G_{<2}|
 =\sqrt{A_0^2+A_1^2+2A_0A_1\cos s}.
\end{equation}
Squaring the displayed vector proves the equality, and projection onto
$-e_2$ followed by the triangle inequality proves
\[
 A_1+A_0\cos s\leq\sigma _1^2\leq A_1+A_0,
 \qquad A_1\leq\sigma _1^2.
\]
The first holding law is $\dot\alpha=F_1=0$.
Consequently $M,G,\zeta,\Theta$ remain at their endpoint values and
$\Omega=0$. The exact linear propagator for arbitrary amplitude data
on this prescribed holding background is
\[
 \begin{pmatrix}\Theta(t)\\\Omega(t)\end{pmatrix}
 =\begin{pmatrix}1&a(t-u)\\0&1\end{pmatrix}
 \begin{pmatrix}\Theta(u)\\\Omega(u)\end{pmatrix}.
\]
For $y=(\Theta,\sigma _0\Omega/\Lambda)$ the upper-right entry is
$\Gamma(t-u)$. Only the selected reference pair has second component
zero; the propagator is not the identity on arbitrary data.

For completeness the next growth interval is explicitly solvable on
this first holding background. Take a fixed $s_2\in(0,1/8]$,
$\Lambda _2>0$, and $p_2=e(s_2)$. Because the first vorticity is zero,
$D_{<2}=0$ and the next phase gradient remains $p_2$. The exact next
growth coefficients and rate are
\[
 a_2=\frac{A_1\sin s_2+A_0\sin(s+s_2)}{\Lambda _2},\quad
 b_2=\Lambda _2\sin s_2,\quad
 \gamma _2=\sqrt{[A_1\sin s_2+A_0\sin(s+s_2)]\sin s_2}>0.
\]
Indeed $-Je(s_2)\cdot G_{<2}$ is the numerator in $a_2$, and
laboratory buoyancy gives $b_2$. An eigenline datum with ratio
$\Omega_2/\Theta_2=\sqrt{b_2/a_2}$ is multiplied by
$e^{\gamma _2(t-t_b)}$, by direct differentiation. If its required
logarithmic gain is $L_2>0$, the first holding integral in the paper's
weighted coordinates is therefore exactly
\[
 H_1=\int_{t_b}^{t_{a2}}\Gamma\,dt=\frac{\Gamma L_2}{\gamma _2}.
\]
Under the actual schedule $s_2=L_2\sigma _0/\sigma _1$, write
$\beta=\arctan[A_0\sin s/(A_1+A_0\cos s)]\in(0,s)$.
Then $G_{<2}=-\sigma _1^2e(-\beta)$, and
$\gamma _2=\sigma _1\sqrt{\sin s_2\sin(s_2+\beta)}$.
Since $s_2+\beta\leq1/4<\pi/2$, this proves the quantitative bounds
\[
 0<H_1=\frac{\sin s\,s_2}
 {\sqrt{\sin s_2\sin(s_2+\beta)}}
 \leq\sin s\frac{s_2}{\sin s_2}
 \leq\frac{\sin s}{1-s_2^2/6}.
\]
The last inequality follows by integrating
$\cos t\geq1-t^2/2$ from zero to $s_2$.
The next transition also has $F_1=F_2=0$: the only term in $F_2$
contains $\Omega_1=0$ and $\zeta_{1,1}=0$.
The displayed constant first-layer state and constant second-stage
geometry solve that closed ODE, whose coefficients are smooth for the
positive vertical components present here. Local uniqueness therefore
extends the constant first-layer state through the second transition.
This proves the first-layer assertion on p.~29 directly.

\paragraph{Exact diffusion check and its feedback consequence.}
For the sine profile and equal physical coefficients
$\kappa_{\rm phys}=\nu_{\rm phys}=d\geq0$, the first affine wave
on its actual rotating base gradient obeys
\[
 \dot\Theta=a\Omega-d\Lambda^2\Theta,\qquad
 \dot\Omega=b\Theta-d\Lambda^2\Omega.
\]
The base remains the displayed evolving affine $G$ because its
Laplacian is zero. Let $(\Theta_0,\Omega_0)$ be the exact inviscid
first-stage solution with the same geometry and data at $t_{\rm in}$.
Then
\[
 (\Theta_d,\Omega_d)(t)
 =e^{-d\Lambda^2(t-t_{\rm in})}(\Theta_0,\Omega_0)(t)
\]
solves both viscous equations by the product rule, with identical
initial data; linear uniqueness proves equality. In particular its
selected vorticity zero and all-trial temperature signs are preserved,
but its gain from the seed is reduced by the exact cost
$d\Lambda^2(L_1+2+C^{-1})/\Gamma$.
This statement uses the fixed inviscid stage times. Its damped
temperature need not cross the paper's threshold at the inviscid time.

During unforced viscous holding, $b=0$ and $\Omega_d=0$, so
\[
 \Theta_d(t)=-P_{d,f}e^{-d\Lambda^2(t-t_b)},\qquad
 G_{<2,d}(t)=(A_0\sin s,
 -A_0\cos s-\Lambda P_{d,f}e^{-d\Lambda^2(t-t_b)}).
\]
Here $P_{d,f}=e^{-d\Lambda^2(t_b-t_{\rm in})}P_f$.
Consequently the next growth coefficient is time dependent:
\[
 a_{2,d}(t)=\frac{\Lambda P_{d,f}e^{-d\Lambda^2(t-t_b)}\sin s_2
                    +A_0\sin(s+s_2)}{\Lambda _2},\qquad
 b_{2,d}=\Lambda _2\sin s_2,
\]
and, before its own diffusion terms are added,
\[
 \dot a_{2,d}(t)=
 -\frac{d\Lambda^3P_{d,f}}{\Lambda _2}
 e^{-d\Lambda^2(t-t_b)}\sin s_2.
\]
Differentiating the prior explicit formulas proves each equality.
Thus the inviscid frozen-background growth rate and holding integral
cannot be carried over unchanged to an unforced diffusive next stage.
Compensating the first-layer scalar diffusion by the exact source
$d\Lambda^2\Theta_{\rm hold}\sin(\Lambda x_2)$, with the constant
$\Theta_{\rm hold}=-P_{d,f}$, makes its scalar amplitude stationary.
This source is an affine-wave source with no
claimed compact support. These identities locate and quantify the
first background feedback term; they assert neither success nor
failure of an infinite viscous iteration.

\paragraph{Exact continuation through the next diffusive growth interval.}
The preceding feedback has an explicit solution beyond merely naming
its coefficients. Take $d>0$, put
\[
 q=d\Lambda^2,\quad \delta_2=d\Lambda _2^2,\quad
 c_0=A_0\sin(s+s_2)\sin s_2>0,\quad
 c_1=\Lambda P_{d,f}\sin^2s_2>0,
\]
and use $r=t-t_b\geq0$. The actual second amplitude pair on the
constant-rotation, diffusing-parent background is
\[
 \Theta_2'=a_{2,d}(r)\Omega_2-\delta_2\Theta_2,\qquad
 \Omega_2'=b_2\Theta_2-\delta_2\Omega_2,
 \quad b_2=\Lambda _2\sin s_2.
\]
Set $W(r)=e^{\delta_2r}\Omega_2(r)$. Since $b_2$ is a positive
constant, differentiating its equation and substituting the first one
gives exactly
\begin{equation}\label{eq:feedbackscalar}
 W''=(c_0+c_1e^{-qr})W,\qquad
 \Omega_2=e^{-\delta_2r}W,\qquad
 \Theta_2=e^{-\delta_2r}W'/b_2.
\end{equation}
Conversely, differentiating these last two expressions verifies the
original pair for every solution of the scalar equation. Thus this is
an invertible morphism between the solution spaces, with initial data
$W(0)=\Omega_2(0)$, $W'(0)=b_2\Theta_2(0)$.

Let $x(r)=x_0e^{-qr/2}$, $x_0=2\sqrt{c_1}/q$, and
$\vartheta=2\sqrt{c_0}/q>0$. Since
$x'=-qx/2$ and $x''=q^2x/4$, the scalar equation is exactly
\[
 x^2\mathcal W_{xx}+x\mathcal W_x-(x^2+\vartheta^2)\mathcal W=0,
 \qquad W(r)=\mathcal W(x(r)).
\]
Here $\mathcal W$ is an explicitly related coordinate representation,
not a replacement of the physical amplitudes.
An exact fundamental pair on $x>0$ is
\[
 f(x)=x^{\vartheta}\sum_{n=0}^\infty a_nx^{2n},\quad
 a_0=1,\quad a_n=\frac{a_{n-1}}{4n(n+\vartheta)},\qquad
 g(x)=f(x)\int_{x_0}^x\frac{d\xi}{\xi f(\xi)^2}.
\]
For each compact interval of $x>0$, the ratio
$|a_nx^{2n}/(a_{n-1}x^{2n-2})|=x^2/[4n(n+\vartheta)]$
tends uniformly to zero. Multiplying terms by any fixed polynomial in
$n$ preserves this fact, so the series and each of its derivatives
converge uniformly there. Substitution is therefore valid term by term;
the coefficient equation is precisely the displayed recurrence.
Every series coefficient is positive, so $f>0$ on $x>0$.
The integral defining $g$ is well defined on compact subintervals.
Differentiating it gives $fg'-f'g=1/x$, and differentiating again
shows $g''+g'/x-(1+\vartheta^2/x^2)g=0$ because $f$ satisfies
the same equation. Hence the Wronskian never vanishes.
Linear uniqueness shows that every solution on $x>0$ is $c_ff+c_gg$.

For actual initial data $W_0=\Omega_2(0)$ and
$W_1=b_2\Theta_2(0)$ the coefficients are exactly
\[
 c_f=\frac{W_0}{f(x_0)},\qquad
 c_g=-\frac{2f(x_0)}q\left[W_1+\frac{qx_0}{2}c_ff'(x_0)\right].
\]
Indeed $g(x_0)=0$, $g'(x_0)=1/[x_0f(x_0)]$, and
$x'(0)=-qx_0/2$ give these two equations by substitution.
Together with \eqref{eq:feedbackscalar}, the formula supplies the
complete physical amplitudes, with no singularity at any finite $r$.
For $d=0$ the original equation is instead
$W''=(c_0+c_1)W$ and its elementary exponential solutions apply
directly; the expressions dividing by $q$ are not used.

The continuation also has rigorous gain bounds. Write
$\gamma_0=\sqrt{c_0}$ and $\gamma_i=\sqrt{c_0+c_1}$, and take
the actual initial growing eigenline
$\Theta_2(0)<0$, $\Omega_2(0)=\sqrt{b_2/a_{2,d}(0)}\Theta_2(0)<0$.
Then $Y=-W$ has $Y(0)=Y_0>0$ and $Y'(0)=\gamma_iY_0>0$.
Its positive coefficient equation excludes a first loss of positivity:
as long as $Y>0$, $Y''>0$ and $Y'\geq Y'(0)>0$.
For all $r\geq0$ set
\[
 Y_-(r)=Y_0\left[\cosh(\gamma_0r)
                  +\frac{\gamma_i}{\gamma_0}\sinh(\gamma_0r)\right],
 \qquad Y_+(r)=Y_0e^{\gamma_ir}.
\]
The difference $Y-Y_-$ has zero initial value and derivative and obeys
$(Y-Y_-)''-\gamma_0^2(Y-Y_-)=c_1e^{-qr}Y\geq0$.
The identity
\[
 Y(r)-Y_-(r)=\int_0^r
 \frac{\sinh(\gamma_0(r-u))}{\gamma_0}
 c_1e^{-qu}Y(u)\,du
\]
is verified by two differentiations and uniqueness. It and its first
derivative are nonnegative. Similarly $Y_+-Y$ has the representation
with kernel $\sinh(\gamma_i(r-u))/\gamma_i$ and source
$c_1(1-e^{-qu})Y(u)\geq0$. Thus
$Y_-\leq Y\leq Y_+$ and $Y_-'\leq Y'\leq Y_+'$.
The exact physical temperature gain consequently satisfies
\[
 e^{-\delta_2r}\left[\cosh(\gamma_0r)
       +\frac{\gamma_0}{\gamma_i}\sinh(\gamma_0r)\right]
 \leq\frac{-\Theta_2(r)}{-\Theta_2(0)}
 \leq e^{(\gamma_i-\delta_2)r}.
\]
The signs of both amplitudes remain negative for finite $r$.
These are finite-interval statements for the exact growth control;
they do not impose a threshold crossing when the diffusion may prevent
one, nor do they claim a subsequent steering or infinite iteration.

\end{document}
